\section{Conclusões}

% Este artigo apresentou o planejamento de uma investigação sobre o impacto de diferentes estratégias de controle populacional em
% algoritmos genéticos aplicados ao problema das \textit{n}-rainhas. A fundamentação teórica revisou conceitos essenciais da área e
% justificou a escolha do problema como cenário experimental. A metodologia detalhou a configuração do algoritmo, os parâmetros
% avaliados e as métricas utilizadas para mensurar desempenho, garantindo reprodutibilidade e rigor na coleta de dados.

% Os experimentos em andamento demonstram que a abordagem proposta fornece uma base sólida para comparar mecanismos de
% substituição populacional. Espera-se que a conclusão da análise permita estabelecer diretrizes práticas sobre como selecionar
% tamanhos de população, taxas de elitismo e proporções de substituição adequadas a diferentes níveis de dificuldade do problema.
% Como trabalhos futuros, pretende-se ampliar o escopo para incluir outras heurísticas evolutivas e investigar técnicas híbridas
% que combinem algoritmos genéticos com métodos exatos ou metaheurísticas complementares.
